\documentclass[11pt,letterpaper]{article}
\usepackage[margin=1in]{geometry}
\usepackage[T1]{fontenc}
\usepackage{lmodern}
\usepackage{amsmath,amssymb,amsthm}
\usepackage{microtype}
\usepackage[hidelinks]{hyperref}
\setlength{\parindent}{0pt}
\setlength{\parskip}{0.48em}
\setlength{\emergencystretch}{1.5em}
\newtheorem{theorem}{Theorem}
\newtheorem{lemma}[theorem]{Lemma}
\newtheorem{corollary}[theorem]{Corollary}
\newcommand{\N}{\mathbb N}
\newcommand{\xM}[1]{Y^{#1}}
\newcommand{\dd}{\partial}
\newcommand{\ee}[2]{\{#1,#2\}}
\newcommand{\code}[1]{\texttt{\small\detokenize{#1}}}
\newcommand{\decl}[1]{{\small\nolinkurl{#1}}}
\pagestyle{plain}

\begin{document}
\begin{center}
{\LARGE Quadratic Markov bases for complete graphs}\\[0.45em]
{\large A disproof of Conjecture 00000000545}\\[0.55em]
C0ldSmi1e \quad $\cdot$ \quad 5 October 2026
\end{center}

\begin{abstract}
Under the standard edge-ring interpretation of graph ring, the proposed
cutoff at five vertices is false. We prove that the toric ideal of every
finite complete graph is generated by a finite set of nonzero quadratic
binomials. The proof connects every degree fiber by replacements of two
edges, with unrestricted edge multiplicities, and then establishes generation
of the full polynomial kernel. The argument works over every nontrivial
commutative ring, hence over every field. In particular, $K_6$ is a
counterexample. The accompanying Lean 4 project checks the entire argument.
\end{abstract}

\section{The statement and the standard objects}
The exact English source is:
\begin{quote}
Definition: Markov bases (quadratic-quadratic). Conjecture: The toric ideal
of the graph ring of $K_n$ has an entirely quadratic Markov basis if and
only if $n\leq 5$ (a criterion beyond Hosten--Sullivant).
\end{quote}
The Chinese version states the same equivalence. Both original texts are
preserved in \code{ORIGINAL.md}. We interpret graph ring as the usual edge
ring, with one generator $x_i x_j$ for each unordered pair of distinct
vertices. No assertion about a different construction, such as a cut ring,
is intended.

Let $V$ be a finite set and let
$E=\{\ee{i}{j}:i,j\in V,\ i\ne j\}$, the edge set of the complete graph.
Let $k$ be a nontrivial commutative ring. Define
\[
 P=k[Y_e:e\in E],\qquad R=k[x_v:v\in V],\qquad
 \varphi:P\longrightarrow R,\quad Y_{\ee{i}{j}}\longmapsto x_i x_j,
\]
with $\varphi$ acting identically on coefficients. The graph ring is
$\operatorname{im}\varphi=k[x_i x_j:i\ne j]$, and its toric ideal is
$I=\ker\varphi$. Over a field, a finite binomial generating set of $I$ is
the algebraic definition of a Markov basis; we also prove its degree-fiber
connectivity directly.

An edge multiset $M$ is an arbitrary vector in $\N^E$. Write $|M|$ for
its number of edges, counting multiplicity, and $\xM{M}=\prod_eY_e^{M(e)}$.
The endpoint multiset $\dd M\in\N^V$ records the number of times each
\emph{labeled vertex} occurs. Thus
\[
 \varphi(\xM{M})=x^{\dd M},\qquad
 |\dd M|=2|M|,\qquad (\dd M)(v)\leq |M|.
\]
The final inequality uses the absence of loops. Equality $\dd M=\dd N$
means equal degree at each labeled vertex, not merely an unordered list of
equal numerical degrees.

\newpage
\section{Inserting an edge with quadratic switches}
A quadratic replacement substitutes a two-edge multiset $A$ by another
two-edge multiset $B$ with $\dd A=\dd B$, while retaining every other edge
occurrence. Repeated edges are allowed. A replacement with $A=B$ is simply
reflexive and will not be included as a generator.

\begin{lemma}[Insertion]
Suppose $a\ne b$ and $\dd M=\dd(\ee{a}{b}+Q)$. There is an edge multiset
$T'$ such that at most two quadratic replacements transform $M$ into
$\ee{a}{b}+T'$.
\end{lemma}
\begin{proof}
The vertex $a$ occurs in $\dd M$, so remove one occurrence of an edge
$\ee{a}{c}$ from $M$. If $c=b$, the desired edge is already present.
Otherwise $b$ is on neither endpoint of the removed edge, since $a\ne b$
and $c\ne b$. The equality of endpoint multisets therefore ensures that
an occurrence of an edge $\ee{b}{d}$ remains. Remove it too, leaving $S$:
\[
 M=\ee{a}{c}+\ee{b}{d}+S.
\]
These are removals of occurrences, so the assertion remains valid with
arbitrary multiplicities.

If $c\ne d$, make the replacement
\[
 \ee{a}{c}+\ee{b}{d}\ \longmapsto\ \ee{a}{b}+\ee{c}{d}.
\]
Both new edges are loopless, and the endpoint multiset is unchanged.

It remains to consider $c=d$. In this case $c\ne a$ and $c\ne b$, because
the two removed edges are loopless. Let $s_c=(\dd S)(c)$ and
$q_c=(\dd Q)(c)$. Comparing the count of the \emph{specific} vertex $c$
and then the total endpoint counts gives
\[
 2+s_c=q_c,\qquad |M|=|Q|+1,\qquad |Q|=|S|+1.
\]
Since $q_c\leq |Q|$, we have
$2+s_c\leq |S|+1$, and hence $s_c<|S|$. Consequently some edge of $S$
avoids $c$: if all its edges contained $c$, each would contribute exactly
one to $s_c$, yielding $s_c=|S|$.

Remove an occurrence of such an edge $\ee{x}{y}$ and write
$S=\ee{x}{y}+T$. Interchange the names $x,y$ if necessary so that $x\ne a$;
this is possible because $x\ne y$. We have $c\ne x$ and $c\ne y$.
The following two replacements, carrying $T$ along, are therefore valid:
\begin{align*}
 \ee{a}{c}+\ee{x}{y}&\ \longmapsto\ \ee{a}{x}+\ee{c}{y},\\
 \ee{a}{x}+\ee{b}{c}&\ \longmapsto\ \ee{a}{b}+\ee{x}{c}.
\end{align*}
Every displayed edge is loopless: the additional inequalities needed are
$a\ne x$, $c\ne y$, $a\ne b$, and $x\ne c$, all already established.
The final multiset is
$\ee{a}{b}+\ee{x}{c}+\ee{c}{y}+T$, as required.
If an intermediate replacement leaves its two-edge multiset unchanged,
it is a zero-step move. Such a case causes no problem and adds no zero
binomial to the eventual basis.
\end{proof}

\newpage
\section{All fibers, with no degree bound}
\begin{theorem}[Quadratic connectivity]
For any edge multisets $M,N$, equality $\dd M=\dd N$ implies that a finite
sequence of quadratic replacements transforms $M$ into $N$.
\end{theorem}
\begin{proof}
Induct on $|N|$, with $M$ arbitrary. If $N$ is empty, then
$2|M|=|\dd M|=|\dd N|=0$, so $M$ is empty as well.

Otherwise remove one occurrence of an edge $e=\ee{a}{b}$ from $N$ and
write $N=e+Q$. The insertion lemma transforms $M$ into $e+T'$.
Every replacement preserves its endpoint multiset, so
\[
 \dd e+\dd T'=\dd M=\dd N=\dd e+\dd Q.
\]
Cancel the common endpoint multiset $\dd e$ to obtain
$\dd T'=\dd Q$. The induction hypothesis transforms $T'$ into $Q$.
Carry the common edge occurrence $e$ along during each of these moves.
Combining the two finite sequences transforms $M$ into $N$.
\end{proof}

This induction concerns all nonnegative exponent vectors. In particular,
it does not enumerate a bounded collection of monomials or rely on a
test of finitely many degree fibers.

Define the following set of polynomials:
\[
 \mathcal B=\left\{\xM{A}-\xM{B}:
 |A|=|B|=2,\ \dd A=\dd B,\ A\ne B\right\},
 \qquad J=(\mathcal B)\subseteq P.
\]
For finite $E$, there are finitely many multisets of two edges; hence
$\mathcal B$ is finite. Every member maps to zero under $\varphi$, so
$J\subseteq I$. The condition $A\ne B$ explicitly excludes zero binomials:
distinct exponent vectors label distinct monomials, and $1\ne0$ in $k$.
Both monomials in each member have total degree two. Since their
coefficients cannot cancel across distinct exponents, each such polynomial
is nonzero, homogeneous of degree two, and has total degree exactly two.

\begin{lemma}[Binomials of arbitrary degree]
If $\dd M=\dd N$, then $\xM{M}-\xM{N}\in J$.
\end{lemma}
\begin{proof}
For a single quadratic replacement $A+C\mapsto B+C$, the corresponding
difference is
\[
 \xM{A+C}-\xM{B+C}=\xM{C}(\xM{A}-\xM{B})\in J.
\]
If $A=B$, this difference is zero, which also belongs to $J$.
For a finite sequence $M_0,\ldots,M_r$ with $M_0=M$ and $M_r=N$, sum
the differences:
\[
 \xM{M}-\xM{N}
   =\sum_{i=0}^{r-1}\bigl(\xM{M_i}-\xM{M_{i+1}}\bigr)\in J.
\]
The connectivity theorem supplies such a sequence whenever the endpoint
multisets agree.
\end{proof}

\newpage
\section{The entire polynomial kernel}
To pass from binomials to all kernel polynomials, the following general
argument is included explicitly. It needs no restriction on polynomial
degree and no assertion that the polynomial being considered is itself
a binomial.

\begin{lemma}[Monomial-map kernel]
Let $A:\N^{(E)}\to\N^{(V)}$ be an additive map, and let
$\Phi:k[Y_e:e\in E]\to k[x_v:v\in V]$ send $Y^u$ to $x^{A(u)}$ and
fix coefficients. Suppose an ideal $L$ contains $Y^u-Y^w$ whenever
$A(u)=A(w)$. Then $\ker\Phi\subseteq L$.
\end{lemma}
\begin{proof}
Let $\pi:P\to P/L$ be the quotient map. The hypothesis says that the
class $\pi(Y^u)$ depends only on $A(u)$. For each realized image exponent
$t$, denote that common class by $g(t)$; set $g(t)=0$ for an exponent not
realized. Equivalently, one may choose one representative monomial for
each realized exponent. This choice does not affect any value used below.

Write an arbitrary polynomial as a finite sum
$p=\sum_u c_uY^u$. Grouping its image by exponent gives
\[
 \Phi(p)=\sum_t\left(\sum_{u:A(u)=t}c_u\right)x^t.
\]
Each displayed sum is finite, with only exponents from the support of $p$
contributing. If $p\in\ker\Phi$, uniqueness of polynomial coefficients
implies
\[
 \sum_{u:A(u)=t}c_u=0\quad\text{for every }t.
\]
This coefficientwise statement is valid over any commutative ring, even
with zero divisors. Consequently
\begin{align*}
 \pi(p)
 &=\sum_u\pi(c_u)\,\pi(Y^u)
   =\sum_u\pi(c_u)\,g(A(u))\\
 &=\sum_t\pi\!\left(\sum_{u:A(u)=t}c_u\right)g(t)=0.
\end{align*}
Thus $p\in L$, proving the inclusion.
\end{proof}

Apply this lemma with $A(M)=\dd M$, $\Phi=\varphi$, and $L=J$.
The preceding binomial lemma verifies its hypothesis for \emph{all}
exponent vectors. Hence $I\subseteq J$. Together with the already proved
inclusion $J\subseteq I$, this yields $I=J=(\mathcal B)$.

\begin{theorem}[The quadratic basis]
For every finite complete graph and every nontrivial commutative ring $k$,
the kernel of its edge-ring parametrization has a finite generating set
$\mathcal B$ consisting entirely of nonzero quadratic pure binomials.
In particular, over every field, $\mathcal B$ is an entirely quadratic
Markov basis.
\end{theorem}
\begin{proof}
Use the set $\mathcal B$ defined above. Its finiteness, the degree and
nonvanishing of its members, and the equality $I=(\mathcal B)$ have all
been established. Quadratic connectivity proves the corresponding
combinatorial Markov property as well.
\end{proof}
No minimality of this generating set is claimed or needed. Including
redundant nonzero generators, or both signs of one binomial, is permitted
in a Markov basis.

\newpage
\section{The counterexample and formal correspondence}
\begin{corollary}
Conjecture 00000000545 is false: $K_6$ has an entirely quadratic Markov
basis, although $6\not\leq5$.
\end{corollary}
\begin{proof}
Take $V=\{0,1,2,3,4,5\}$ in the quadratic-basis theorem. The implication
``has an entirely quadratic Markov basis $\Rightarrow n\leq5$'' fails.
Since $6>0$, the counterexample also applies when the conjecture is
quantified only over positive integers.
\end{proof}

The formal proof uses Lean 4.19.0 and mathlib at commit
\begin{center}\small\ttfamily
c44e0c8ee63ca166450922a373c7409c5d26b00b.
\end{center}
The package pins all dependencies. Its polynomial rings are mathlib's
\decl{MvPolynomial}; the toric ideal is the actual \decl{RingHom.ker}.
The correspondence between arbitrary edge multisets and natural exponent
vectors is mathlib's additive equivalence \decl{Multiset.toFinsupp}.
The property \decl{CompleteGraph.HasQuadraticMarkovBasis} explicitly
requires a finite binomial set spanning this kernel, distinct exponent
vectors, degree two for both monomials, and polynomial total degree two.

The principal declarations are as follows; module names indicate files,
while every displayed qualified theorem name is its actual Lean name.
\begin{itemize}
\raggedright
\setlength{\itemsep}{0.35em}
\setlength{\parskip}{0pt}
\item \code{Connectivity.lean}: edges are all cardinality-two vertex
finsets. \decl{CompleteGraph.insert_edge} proves the insertion lemma, and
\decl{CompleteGraph.connected_of_degree_eq} proves unrestricted fiber
connectivity, including repeated edge occurrences.
\item \code{IdealBridge.lean}:
\decl{IdealBridge.ker_monomialMap_le_of_fiber_binomials} proves the
arbitrary-polynomial kernel lemma. The substitution identity is
\texttt{\small IdealBridge.toricMap\_eq\_eval}${}_2$\texttt{\small Hom}.
\item \code{Conjecture545.lean}:
\decl{CompleteGraph.graphRingMap_edge} identifies the parametrization as
$Y_{\ee{i}{j}}\mapsto x_i x_j$.
\decl{CompleteGraph.toricIdeal_eq_span_quadraticBinomials} proves the
full ideal equality. The declarations
\decl{CompleteGraph.quadraticBinomials_finite} and
\decl{CompleteGraph.quadraticBinomials_totalDegree} verify finiteness
and exact degree two.
\item \decl{CompleteGraph.hasQuadraticMarkovBasis} gives the general
existence theorem, and \decl{CompleteGraph.counterexample_six} supplies
$K_6$. The final positive-domain negation is
\decl{CompleteGraph.conjecture545_false_positive}; the unrestricted
natural-number version is \decl{CompleteGraph.conjecture545_false}.
\end{itemize}

The full project build and separate warnings-as-errors replays of all five
Lean files pass. The compiled inventory covers all 93 authored and
generated mathematical declarations, with axiom dependencies contained in
Lean's standard \decl{propext}, \decl{Quot.sound}, and
\decl{Classical.choice}. No additional axiom, unproved placeholder, or
native-decision shortcut is used. All nine pinned dependency revisions and
their tracked source trees were checked. There is no auxiliary numerical
computation on which the proof relies. \code{Audit.lean},
\code{Check.lean}, the declaration inventory, and the verification record
accompany the project; they supplement, rather than replace, the proof.

\end{document}
